\subsection{Exact Periodic Gap in the One-Species, Zero-Phase Model}
\label{sec:pvbs_periodic_gap}

For positive real hopping $\mu$, the one-particle dispersion in
\cite[Section~II, after equation~(8)]{BachmannNachtergaele2012PVBS} is
$1-2\mu\cos(k)/(1+\mu^2)$ at zero phase. The general multi-species
thermodynamic gap formula stated there is a conjecture. The finite-ring
result below is a proved one-species specialization, with an explicit
operator argument; its scope is recorded
in~\cite{gap:bn12_pvbs_periodic_gap}.

\begin{lemma}[Cyclic site and window permutations]
    \label{lem:pvbs_cyclic_geometry}
    \lean{MPSTensor.cyclicSuccessorEquiv,
        MPSTensor.cyclicSuccessorEquiv_apply,
        MPSTensor.cyclicSuccessorEquiv_ne,
        MPSTensor.finEquiv_cyclicSuccessorEquiv,
        MPSTensor.finEquiv_cyclicSuccessorEquiv_symm,
        MPSTensor.cyclicCfg_extractWindow_comp_swap_of_le}
    \leanok
    On a ring of at least two sites, the cyclic successor is a permutation
    without fixed points. In cyclic-group coordinates it is addition by
    one, and its inverse is subtraction by one. Exchanging two entries
    of an extracted window exchanges the corresponding physical sites
    when that window is inserted back into the configuration.
\end{lemma}
\begin{proof}
    \leanok
    Transport addition through the equivalence between finite site indices
    and the cyclic group. For the exchange identity, express each
    transposition as two function updates and commute those updates with
    insertion of the cyclic window.
\end{proof}

\begin{lemma}[Normalized local PVBS projector]
    \label{lem:pvbs_local_projector}
    \lean{MPSTensor.parentInteraction_pvbs_two_apply}
    \leanok
    \uses{thm:pvbs_open_parent}
    At real hopping $\mu$, the canonical two-site parent interaction is
    \begin{align}
        h_\mu & =|11\rangle\langle11|+
        \frac{(|01\rangle-\mu|10\rangle)
            (\langle01|-\mu\langle10|)}{1+\mu^2}.
        \label{eq:pvbs_normalized_interaction}
    \end{align}
    Equivalently, with occupation projections $n_L,n_R$, it is
    $(\mu^2n_L+n_R-\mu T)/(1+\mu^2)$, where $T$ exchanges $01$ and $10$
    and annihilates $00$ and $11$.
\end{lemma}
\begin{proof}
    \leanok
    The image of the displayed operator is orthogonal to the two-site
    vacuum-particle space, and subtracting its image from any vector leaves
    a vector in that space. Uniqueness of orthogonal projection identifies
    the displayed operator with the actual parent interaction.
\end{proof}

\begin{theorem}[Periodic occupation decomposition]
    \label{thm:pvbs_periodic_occupation}
    \lean{MPSTensor.localTermES_pvbs_two_apply,
        MPSTensor.localTermES_pvbs_decomposition,
        MPSTensor.parentHamiltonianES_pvbs_decomposition}
    \leanok
    \uses{lem:pvbs_local_projector,lem:pvbs_cyclic_geometry}
    On every ring of $N\ge2$ sites, the actual PVBS parent Hamiltonian satisfies
    \begin{align}
        H_\mu & =\frac{2\mu}{1+\mu^2}H_1+
        \frac{(\mu-1)^2}{1+\mu^2}\mathcal N,
              & \mathcal N|\sigma\rangle  & =
        \left(\sum_i\sigma_i\right)|\sigma\rangle.
        \label{eq:pvbs_periodic_number}
    \end{align}
    The sum defining $H_\mu$ contains both oriented windows when $N=2$.
\end{theorem}
\begin{proof}
    \leanok
    The difference between $h_\mu$ and $2\mu h_1/(1+\mu^2)$ is the
    diagonal field $[(\mu^2-\mu)n_L+(1-\mu)n_R]/(1+\mu^2)$.
    Summing cyclically gives the same total occupation at the left and
    right ends of the windows, yielding (\ref{eq:pvbs_periodic_number}).
\end{proof}

\begin{theorem}[Uniform periodic PVBS lower gap bound]
    \label{thm:pvbs_uniform_lower_gap}
    \lean{MPSTensor.parentHamiltonianES_pvbs_energy_decomposition,
        MPSTensor.pvbs_particle_number_ge_norm_sq,
        MPSTensor.parentHamiltonianES_pvbs_energy_lower_bound,
        MPSTensor.pvbs_zero_of_mem_orthogonal_parent_kernel,
        MPSTensor.parentHamiltonianES_pvbs_gap}
    \leanok
    \uses{thm:pvbs_periodic_occupation,thm:pvbs_periodic_parent}
    For $\mu\ge0$, set $c(\mu)=(\mu-1)^2/(1+\mu^2)$.
    If $v$ has zero vacuum amplitude, then
    $\langle v,H_\mu v\rangle\ge c(\mu)\|v\|^2$.
    Consequently, for every $N\ge2$ and $v\perp\ker H_\mu$,
    \begin{align}
        c(\mu)\|v\| & \le\|H_\mu v\|.
        \label{eq:pvbs_uniform_norm_gap}
    \end{align}
\end{theorem}
\begin{proof}
    \leanok
    The critical parent $H_1$ is positive. Every nonvacuum configuration
    contains at least one particle, so $\langle v,\mathcal N v\rangle
        \ge\|v\|^2$ when the vacuum coefficient vanishes. Apply
    (\ref{eq:pvbs_periodic_number}). The vacuum belongs to the actual
    kernel, so kernel-orthogonal vectors meet this condition. The norm
    bound follows from Cauchy--Schwarz.
\end{proof}

\begin{theorem}[Attainment and exact noncritical finite-ring gap]
    \label{thm:pvbs_exact_periodic_gap}
    \lean{MPSTensor.parentHamiltonianES_pvbs_critical_uniformParticle,
        MPSTensor.parentHamiltonianES_pvbs_uniformParticle,
        MPSTensor.parentHamiltonianES_pvbs_gap_le,
        MPSTensor.parentHamiltonianES_pvbs_gap_exact}
    \leanok
    \uses{thm:pvbs_uniform_lower_gap}
    The uniform one-particle vector $W_N=\sum_j|0\cdots1_j\cdots0\rangle$
    satisfies $H_\mu W_N=c(\mu)W_N$ for every real $\mu$ and $N\ge2$.
    For $\mu>0$ and $\mu\ne1$, $c(\mu)>0$ is the largest possible
    norm-gap constant in (\ref{eq:pvbs_uniform_norm_gap}), hence the exact
    gap above the unique periodic vacuum.
    At $\mu=1$, the W vector is instead another zero-energy state;
    this does not assert that the finite-ring gap above the enlarged
    kernel is zero.
\end{theorem}
\begin{proof}
    \leanok
    The uniform one-particle vector is invariant under exchange of two
    sites and has no adjacent double occupation. Every critical local
    term therefore annihilates it. It has particle number one, so
    (\ref{eq:pvbs_periodic_number}) gives its eigenvalue $c(\mu)$.
    For noncritical hopping, this is positive, and symmetry of the
    Hamiltonian makes the eigenvector orthogonal to its kernel.
    Applying any proposed norm-gap bound to this nonzero vector gives
    the matching upper bound $c(\mu)$.
\end{proof}
